\scanpage{84}
% archive foliation: 85

Exemples de groupes de BT.

a)\quad Soit $G$ un schéma en groupes lisse sur $S$, \deleted{à fibres connexes} extension d'un schéma abélien\inserted{$A$} par un tore\inserted{$T$}. Alors $p.\mathrm{id}_G$ est un épimorphisme, et les $G(n)=\operatorname{Ker}p^n\mathrm{id}_G$ sont des groupes finis localement libres (commencer par cas d'un schéma abélien). Donc $G(\infty)=\lim G(n)$ est un groupe de BT.\inserted{Si $A$ ($T$) de dim relative $a(m)$, $G(\infty)$ de rang $2a+m$.} Cas important: $\mu\deleted{\unsure{(\infty)}}=\underline{G}_m(\infty)$.

b)\quad Supposons que $S$ soit \deleted{à caractéristique résiduelle égale à $p$}\inserted{tel que $p$ soit loc. nilpotent sur $S$.} et soit $G$ un groupe \inserted{de Lie}formel sur $S\deleted{\unsure{x}}$ (cette notion sera précisée plus tard, pour ceux qui ne la connaitraient pas encore \ldots).\inserted{On a alors $G=G(\infty)=\lim G(n)$.} Supposons que $p.\mathrm{id}_G$ soit un épimorphisme; lorsque $S$ est artinien local, comme unique point $s\in S$, l'hypothèse revient à dire que $p.\mathrm{id}_{G_s}$ dans $G_s$ est un épimorphisme, ou encore que $G_{\bar{s}}$ ne contient pas de sous-groupe de Lie formel isomorphe à $\underline{G}_a$. (Mais si $S$ n'est pas réduit à un point, p. ex. si $S$ est un trait, la condition que $p.\mathrm{id}_G$ soit un épimorphisme \uline{n'est pas} une simple condition fibre par fibré. Prendre par exemple le groupe de Lie formel associé à\deleted{un groupe} la courbe ell\deleted{\unsure{e}}iptique modulaire, au voisinage d'un point de Hasse nul ..\inserted{$(\unsure{\dagger|*})$}) Il devrait être vrai que \deleted{les} $G(1)$ dont les $G(n)$ sont finis localement libres (plus généralement, si on a un épimorphisme de groupes de Lie formels de même dimension relative, le noyau devrait être fini localement libre, et réciproquement \ldots). C'est en tous cas le cas si $S$ est artinien; supposons cette condition supplémentaire vérifiée. On trouve alors \inserted{que $G$ est}un groupe de BT \deleted{$G(\infty)$}. De plus, les $G(n)$ sont radiciels, et on \inserted{prouvera par la suite que l'}trouve une équivalence entre la catégorie des groupes de BT $G$ sur $S$ tel que $G(1)$ soit radiciel (ou ce qui revient au même, les $G(n)$ le soient -- on dit alors que $G$ est à fibres connexes), et la catégorie des groupes de Lie formels \inserted{$G$}sur $S$ tels que $p.\mathrm{id}_G$ soit un épimorphisme et tel que $G(1)$ soit fini localement libre (ces deux

\margin{$(\unsure{\dagger|*})$ La condition suppl. à rajouter devrait être : fibres de ``hauteur'' localement constante}
